% Options for packages loaded elsewhere
\PassOptionsToPackage{unicode}{hyperref}
\PassOptionsToPackage{hyphens}{url}
\documentclass[
  french,
]{article}
\usepackage{xcolor}
\usepackage{amsmath,amssymb}
\setcounter{secnumdepth}{-\maxdimen} % remove section numbering
\usepackage{iftex}
\ifPDFTeX
  \usepackage[T1]{fontenc}
  \usepackage[utf8]{inputenc}
  \usepackage{textcomp} % provide euro and other symbols
\else % if luatex or xetex
  \usepackage{unicode-math} % this also loads fontspec
  \defaultfontfeatures{Scale=MatchLowercase}
  \defaultfontfeatures[\rmfamily]{Ligatures=TeX,Scale=1}
\fi
\usepackage{lmodern}
\ifPDFTeX\else
  % xetex/luatex font selection
\fi
% Use upquote if available, for straight quotes in verbatim environments
\IfFileExists{upquote.sty}{\usepackage{upquote}}{}
\IfFileExists{microtype.sty}{% use microtype if available
  \usepackage[]{microtype}
  \UseMicrotypeSet[protrusion]{basicmath} % disable protrusion for tt fonts
}{}
\makeatletter
\@ifundefined{KOMAClassName}{% if non-KOMA class
  \IfFileExists{parskip.sty}{%
    \usepackage{parskip}
  }{% else
    \setlength{\parindent}{0pt}
    \setlength{\parskip}{6pt plus 2pt minus 1pt}}
}{% if KOMA class
  \KOMAoptions{parskip=half}}
\makeatother
\usepackage{color}
\usepackage{fancyvrb}
\newcommand{\VerbBar}{|}
\newcommand{\VERB}{\Verb[commandchars=\\\{\}]}
\DefineVerbatimEnvironment{Highlighting}{Verbatim}{commandchars=\\\{\}}
% Add ',fontsize=\small' for more characters per line
\newenvironment{Shaded}{}{}
\newcommand{\AlertTok}[1]{\textcolor[rgb]{1.00,0.00,0.00}{\textbf{#1}}}
\newcommand{\AnnotationTok}[1]{\textcolor[rgb]{0.38,0.63,0.69}{\textbf{\textit{#1}}}}
\newcommand{\AttributeTok}[1]{\textcolor[rgb]{0.49,0.56,0.16}{#1}}
\newcommand{\BaseNTok}[1]{\textcolor[rgb]{0.25,0.63,0.44}{#1}}
\newcommand{\BuiltInTok}[1]{\textcolor[rgb]{0.00,0.50,0.00}{#1}}
\newcommand{\CharTok}[1]{\textcolor[rgb]{0.25,0.44,0.63}{#1}}
\newcommand{\CommentTok}[1]{\textcolor[rgb]{0.38,0.63,0.69}{\textit{#1}}}
\newcommand{\CommentVarTok}[1]{\textcolor[rgb]{0.38,0.63,0.69}{\textbf{\textit{#1}}}}
\newcommand{\ConstantTok}[1]{\textcolor[rgb]{0.53,0.00,0.00}{#1}}
\newcommand{\ControlFlowTok}[1]{\textcolor[rgb]{0.00,0.44,0.13}{\textbf{#1}}}
\newcommand{\DataTypeTok}[1]{\textcolor[rgb]{0.56,0.13,0.00}{#1}}
\newcommand{\DecValTok}[1]{\textcolor[rgb]{0.25,0.63,0.44}{#1}}
\newcommand{\DocumentationTok}[1]{\textcolor[rgb]{0.73,0.13,0.13}{\textit{#1}}}
\newcommand{\ErrorTok}[1]{\textcolor[rgb]{1.00,0.00,0.00}{\textbf{#1}}}
\newcommand{\ExtensionTok}[1]{#1}
\newcommand{\FloatTok}[1]{\textcolor[rgb]{0.25,0.63,0.44}{#1}}
\newcommand{\FunctionTok}[1]{\textcolor[rgb]{0.02,0.16,0.49}{#1}}
\newcommand{\ImportTok}[1]{\textcolor[rgb]{0.00,0.50,0.00}{\textbf{#1}}}
\newcommand{\InformationTok}[1]{\textcolor[rgb]{0.38,0.63,0.69}{\textbf{\textit{#1}}}}
\newcommand{\KeywordTok}[1]{\textcolor[rgb]{0.00,0.44,0.13}{\textbf{#1}}}
\newcommand{\NormalTok}[1]{#1}
\newcommand{\OperatorTok}[1]{\textcolor[rgb]{0.40,0.40,0.40}{#1}}
\newcommand{\OtherTok}[1]{\textcolor[rgb]{0.00,0.44,0.13}{#1}}
\newcommand{\PreprocessorTok}[1]{\textcolor[rgb]{0.74,0.48,0.00}{#1}}
\newcommand{\RegionMarkerTok}[1]{#1}
\newcommand{\SpecialCharTok}[1]{\textcolor[rgb]{0.25,0.44,0.63}{#1}}
\newcommand{\SpecialStringTok}[1]{\textcolor[rgb]{0.73,0.40,0.53}{#1}}
\newcommand{\StringTok}[1]{\textcolor[rgb]{0.25,0.44,0.63}{#1}}
\newcommand{\VariableTok}[1]{\textcolor[rgb]{0.10,0.09,0.49}{#1}}
\newcommand{\VerbatimStringTok}[1]{\textcolor[rgb]{0.25,0.44,0.63}{#1}}
\newcommand{\WarningTok}[1]{\textcolor[rgb]{0.38,0.63,0.69}{\textbf{\textit{#1}}}}
\usepackage{longtable,booktabs,array}
\newcounter{none} % for unnumbered tables
\usepackage{calc} % for calculating minipage widths
% Correct order of tables after \paragraph or \subparagraph
\usepackage{etoolbox}
\makeatletter
\patchcmd\longtable{\par}{\if@noskipsec\mbox{}\fi\par}{}{}
\makeatother
% Allow footnotes in longtable head/foot
\IfFileExists{footnotehyper.sty}{\usepackage{footnotehyper}}{\usepackage{footnote}}
\makesavenoteenv{longtable}
\ifLuaTeX
\usepackage[bidi=basic,shorthands=off]{babel}
\else
\usepackage[bidi=default,shorthands=off]{babel}
\fi
\ifLuaTeX
  \usepackage{selnolig} % disable illegal ligatures
\fi
\setlength{\emergencystretch}{3em} % prevent overfull lines
\providecommand{\tightlist}{%
  \setlength{\itemsep}{0pt}\setlength{\parskip}{0pt}}
\usepackage{bookmark}
\IfFileExists{xurl.sty}{\usepackage{xurl}}{} % add URL line breaks if available
\urlstyle{same}
\hypersetup{
  pdftitle={Localisation de sources ponctuelles : LASSO{,} BLASSO{,} maximum de vraisemblance et hypothèses de Duval--Peyré},
  pdflang={fr},
  hidelinks,
  pdfcreator={LaTeX via pandoc}}

\title{Localisation de sources ponctuelles : LASSO, BLASSO, maximum de
vraisemblance et hypothèses de Duval--Peyré}
\usepackage{etoolbox}
\makeatletter
\providecommand{\subtitle}[1]{% add subtitle to \maketitle
  \apptocmd{\@title}{\par {\large #1 \par}}{}{}
}
\makeatother
\subtitle{Compte rendu du dossier sparse\_sources}
\author{}
\date{}

\begin{document}
\maketitle

\section{Compte rendu : localisation de sources
ponctuelles}\label{compte-rendu-localisation-de-sources-ponctuelles}

Branche \texttt{travail-sur-le-probleme-inverse}, dossier
\texttt{sparse\_sources/}. Les figures citées sont produites par les
scripts \texttt{experiments.py}, \texttt{experiments\_1r2.py} et
\texttt{experiments\_projet.py} dans \texttt{sparse\_sources/results/}.
Les tableaux chiffrés viennent des fichiers CSV du même dossier.

Les formules sont écrites en LaTeX (\texttt{\$...\$} et
\texttt{\$\$...\$\$}). Le document se lit directement dans un aperçu
Markdown (VS Code, GitHub) et se convertit en LaTeX ou en PDF avec
Pandoc (voir la section 10).

\subsection{0. Résumé}\label{ruxe9sumuxe9}

\begin{enumerate}
\def\labelenumi{\arabic{enumi}.}
\tightlist
\item
  \textbf{Outils implémentés} : un LASSO positif sur grille (FISTA,
  sklearn, ensemble actif), un BLASSO sans grille par Frank-Wolfe
  glissant, une fonction \texttt{check\_hypotheses} qui vérifie la
  condition de source non dégénérée de Duval--Peyré, et une méthode de
  \textbf{maximum de vraisemblance gloutonne} avec choix du nombre de
  sources par BIC (\texttt{greedy\_ml}), accompagnée de la \textbf{borne
  de Cramér--Rao} (\texttt{cramer\_rao}). Le code est couvert par 70
  tests pytest.
\item
  \textbf{Le BLASSO standard échoue toujours avec le noyau du projet}
  \(G(x,z) = 1/\lVert x-z\rVert^2\) (capteurs dans le plan) : la
  condition de Duval--Peyré est fausse dans 100 \% des tirages. La cause
  est mathématique : \(G(x,\cdot)\) est sous-harmonique et
  \(\lVert\varphi\rVert\) varie d'un facteur 215 sur le domaine.
\item
  \textbf{Le BLASSO pondéré} (pénalité
  \(\lambda\sum_k I_k\lVert\varphi(z_k)\rVert^{\alpha}\), avec
  \(\alpha\approx 1\)) corrige ce défaut : 1 ou 2 sources sont
  retrouvées dans 95 à 100 \% des tirages, et \texttt{check\_hypotheses}
  prédit la réussite (94 \% de réussite si la condition est vraie, 2 à 6
  \% sinon).
\item
  \textbf{Dans le scénario réel du projet} (grosse source centrale
  d'intensité 5000 retirée, bruit relatif), la meilleure méthode est le
  \textbf{maximum de vraisemblance glouton + BIC}. Elle bat Prony et le
  BLASSO dans presque toutes les cases, et son erreur est proche de la
  borne de Cramér--Rao, qui est la meilleure précision possible (section
  5.5 ; tableau de choix des méthodes en section 11).
\item
  \textbf{La limite principale est physique, pas algorithmique} : le
  bruit de la source centrale reste dans les mesures après soustraction.
  Il multiplie l'erreur minimale possible par un facteur d'environ 40
  par rapport au même problème sans source centrale (3 sources,
  \(\sigma = 10^{-3}\) : 100 \% de réussite sans source centrale, 40 \%
  avec). Dans la configuration par défaut
  (\(N\sim\mathrm{Poisson}(20)\), \(\sigma = 1\,\%\)), seules 0 à 4
  sources par tirage sont détectables, quelle que soit la méthode.
\end{enumerate}

\subsection{1. Cadre et notations}\label{cadre-et-notations}

\textbf{Domaine et capteurs.} Les sources sont dans
\(D = [0,10]\times[0,20]\) (\texttt{config/simulation.yaml}). Les
\(M = 64\) capteurs sont équirépartis sur le cercle circonscrit à \(D\),
de centre \(c = (5,10)\) et de rayon \(R = \sqrt{125}\approx 11{,}18\),
dans le plan des sources (\texttt{config/sensors.yaml}). Ce cercle passe
par les coins de \(D\) : certains capteurs sont à environ \(0{,}3\) d'un
coin.

\textbf{Sources et mesures.} Les sources forment une mesure positive

\[
s = \sum_{k=1}^{K} I_k\,\delta_{z_k}, \qquad I_k > 0,\quad z_k \in D,
\]

avec \(K\) inconnu. Le noyau du projet
(\texttt{particle\_sources.model.field.intensity\_at}) est

\[
G(x_i, z) = \frac{1}{\lVert x_i - z\rVert^2 + h^2},
\qquad
\varphi(z) = \big(G(x_1,z),\dots,G(x_M,z)\big) \in \mathbb{R}^M .
\]

On a \(h = 0\) dans le projet (paramètre \texttt{eps}). Un capteur placé
à la hauteur \(h\) au-dessus du plan voit exactement ce noyau ; \(h\)
sert donc à étudier des capteurs hors du plan. Les mesures sont

\[
m = A(s) + w, \qquad A(s) = \sum_{k=1}^{K} I_k\,\varphi(z_k).
\]

\textbf{Deux modèles de bruit.}

\begin{itemize}
\tightlist
\item
  Bruit additif (études génériques, sections 5.1 à 5.3) :
  \(w \sim \mathcal{N}(0,\sigma^2 I_M)\), avec
  \(\sigma = \text{niveau}\times\mathrm{rms}(A(s))\).
\item
  Bruit relatif du projet (sections 5.4 et 5.5) :
  \(y_i = y_i^{\mathrm{exact}}\,(1+\sigma\varepsilon_i)\), avec
  \(\varepsilon_i\sim\mathcal{N}(0,1)\). Les mesures brutes \(y\)
  contiennent la source centrale \(I_c\,\varphi(c)\), avec
  \(I_c = 50\times 100 = 5000\). Le projet la retire :
  \(m = y - I_c\,\varphi(c)\). Le bruit, lui, reste, avec un écart type
  par capteur
\end{itemize}

\[
s_i = \sigma\, y_i \approx \sigma\Big(\frac{I_c}{R^2} + A(s)_i\Big) = \sigma\,(40 + A(s)_i).
\]

La source centrale est au centre du cercle des capteurs : elle apporte
exactement \(I_c/R^2 = 40\) à chaque capteur. Le signal utile vaut
environ 7 en moyenne quadratique. Le bruit est donc dominé par la source
centrale.

\subsubsection{1.1 Problèmes
d'optimisation}\label{probluxe8mes-doptimisation}

\textbf{BLASSO positif} (sans grille) :

\[
\min_{s\ge 0}\; \tfrac12\,\lVert m - A(s)\rVert_2^2 + \lambda\,\lVert s\rVert_{\mathrm{TV}},
\qquad \lVert s\rVert_{\mathrm{TV}} = \sum_k I_k .
\]

\textbf{LASSO positif sur grille} (positions restreintes aux nœuds
\(y_1,\dots,y_P\)) :

\[
\min_{c\in\mathbb{R}^P_{+}}\; \tfrac12\,\lVert m - \mathbf{G}c\rVert_2^2 + \lambda\sum_{j=1}^{P} c_j,
\qquad \mathbf{G}_{ij} = G(x_i, y_j)\in\mathbb{R}^{M\times P}.
\]

\textbf{BLASSO pondéré} (ajout par rapport à la demande) :

\[
\min_{K,\,I\ge 0,\,z}\; \tfrac12\,\Big\lVert m - \sum_k I_k\varphi(z_k)\Big\rVert_2^2
+ \lambda\sum_k I_k\,\lVert\varphi(z_k)\rVert^{\alpha}.
\]

\subsubsection{\texorpdfstring{1.2 Ce que représente
\(\alpha\)}{1.2 Ce que représente \textbackslash alpha}}\label{ce-que-repruxe9sente-alpha}

\(\alpha \ge 0\) est \textbf{l'exposant de pondération de la pénalité}.
Le BLASSO pondéré est exactement le BLASSO standard pour le noyau
modifié

\[
\psi(z) = \frac{\varphi(z)}{\lVert\varphi(z)\rVert^{\alpha}},
\qquad \tilde I_k = I_k\,\lVert\varphi(z_k)\rVert^{\alpha},
\]

car \(\sum_k \tilde I_k\,\psi(z_k) = \sum_k I_k\,\varphi(z_k)\) : le
modèle physique et les mesures sont inchangés, seule la pénalité change.

\begin{itemize}
\tightlist
\item
  \(\alpha = 0\) : BLASSO standard ; chaque unité d'intensité coûte
  \(\lambda\), où que soit la source.
\item
  \(\alpha = 1\) : noyau normalisé, \(\lVert\psi(z)\rVert = 1\) partout.
  Une source coûte \(\lambda\) par unité de \textbf{signal} produit, et
  non par unité de masse.
\item
  \(\alpha > 1\) : on favorise encore plus les sources loin des
  capteurs.
\end{itemize}

Pourquoi c'est nécessaire avec le noyau en \(1/r^2\) :
\(\lVert\varphi(z)\rVert\) vaut \(0{,}064\) au centre de \(D\) et
\(13{,}7\) près d'un coin, soit un rapport de 215. Avec \(\alpha = 0\),
une source près d'un capteur produit 215 fois plus de signal par unité
de pénalité ; le BLASSO remplace alors une source centrale par des
sources parasites au bord. Avec \(\alpha = 1\), on montre par
Cauchy--Schwarz qu'une source seule est toujours bien identifiée
(section 5.2). La valeur conseillée est \(\alpha = 1\) (ou \(1{,}25\)).
Dans le code : \texttt{make\_setup(normalize=True)} (\(\alpha = 1\)) ou
\texttt{make\_setup(alpha=...)}.

\subsubsection{1.3 Certificat et
optimalité}\label{certificat-et-optimalituxe9}

Pour le résidu \(r = m - A(s)\), le certificat dual est

\[
\eta(z) = \frac{1}{\lambda}\sum_{i=1}^{M} r_i\,G(x_i, z) = \frac{\langle\varphi(z), r\rangle}{\lambda}.
\]

La mesure \(s\) est solution du BLASSO positif \textbf{si et seulement
si} \(\eta(z)\le 1\) sur \(D\) et \(\eta(z_k) = 1\) en chaque source.
C'est le critère d'arrêt du Frank-Wolfe et le test « certifié » du
LASSO.

\subsubsection{1.4 Théorème de Duval--Peyré (2015) et hypothèses
vérifiées}\label{thuxe9oruxe8me-de-duvalpeyruxe9-2015-et-hypothuxe8ses-vuxe9rifiuxe9es}

On pose

\[
\Gamma = \big[\varphi(z_1),\dots,\varphi(z_K)\;\big\vert\;\partial_1\varphi(z_1),\dots,\partial_1\varphi(z_K)\;\big\vert\;\partial_2\varphi(z_1),\dots,\partial_2\varphi(z_K)\big] \in \mathbb{R}^{M\times 3K},
\]

\[
p_V = (\Gamma^{\top})^{+}\begin{pmatrix}\mathbf{1}_K\\ \mathbf{0}_{2K}\end{pmatrix},
\qquad \eta_V(z) = \langle\varphi(z), p_V\rangle .
\]

Le pré-certificat \(\eta_V\) est le certificat de norme minimale qui
vaut 1 aux sources avec un gradient nul. Les hypothèses sont :

\begin{itemize}
\tightlist
\item
  \textbf{(H1)} \(\operatorname{rang}\Gamma = 3K\) (il faut \(M\ge 3K\))
  ;
\item
  \textbf{(H2)} \(\eta_V(z) < 1\) pour tout
  \(z\in D\setminus\{z_1,\dots,z_K\}\), bord compris ;
\item
  \textbf{(H3)} \(\nabla^2\eta_V(z_k)\prec 0\) pour tout \(k\).
\end{itemize}

\textbf{Théorème.} Si (H1)--(H3) sont vraies et si les \(z_k\) sont
intérieurs à \(D\), il existe \(\lambda_0, c_0 > 0\) tels que, pour
\(\lambda\le\lambda_0\) et \(\lVert w\rVert/\lambda\le c_0\), le BLASSO
a exactement \(K\) Diracs, à distance \(O(\lVert w\rVert)\) des vraies
sources. Le théorème est \textbf{qualitatif} : \(\lambda_0\) et \(c_0\)
ne sont pas calculables. On mesure donc aussi la marge et le
conditionnement de \(\Gamma\) (section 5.3).

\subsection{2. Méthodes
implémentées}\label{muxe9thodes-impluxe9mentuxe9es}

\subsubsection{\texorpdfstring{2.1 Noyau
(\texttt{kernel.py})}{2.1 Noyau (kernel.py)}}\label{noyau-kernel.py}

On réutilise \texttt{InverseSquareKernel} du projet, avec des dérivées
analytiques. En notant \(q = \lVert z-x\rVert^2 + h^2\) :

\[
\nabla_z G = -\frac{2(z-x)}{q^2},
\qquad
\nabla_z^2 G = -\frac{2\,\mathrm{Id}}{q^2} + \frac{8\,(z-x)(z-x)^{\top}}{q^3}.
\]

\begin{itemize}
\tightlist
\item
  \texttt{Setup} regroupe les capteurs, \(h\), \(D\), la pondération
  \(\alpha\) et des \textbf{poids par capteur} \(w_i\) (blanchiment,
  section 2.6) ; \texttt{make\_setup()} lit \texttt{config/}.
\item
  \texttt{WeightedKernel} calcule
  \(\psi = \varphi/\lVert\varphi\rVert^{\alpha}\) avec ses dérivées
  analytiques ; \texttt{WhitenedKernel} calcule \((w_i\,G(x_i,z))_i\).
\item
  Toutes les dérivées sont vérifiées contre des différences finies
  (erreur relative \(\le 10^{-9}\)).
\end{itemize}

\subsubsection{\texorpdfstring{2.2 Briques communes
(\texttt{common.py})}{2.2 Briques communes (common.py)}}\label{briques-communes-common.py}

{\def\LTcaptype{none} % do not increment counter
\begin{longtable}[]{@{}
  >{\raggedright\arraybackslash}p{(\linewidth - 2\tabcolsep) * \real{0.5000}}
  >{\raggedright\arraybackslash}p{(\linewidth - 2\tabcolsep) * \real{0.5000}}@{}}
\toprule\noalign{}
\begin{minipage}[b]{\linewidth}\raggedright
Fonction
\end{minipage} & \begin{minipage}[b]{\linewidth}\raggedright
Rôle
\end{minipage} \\
\midrule\noalign{}
\endhead
\bottomrule\noalign{}
\endlastfoot
\texttt{forward(z,\ I)} & \(A(s) = \sum_k I_k\,\varphi(z_k)\) \\
\texttt{certificate}, \texttt{max\_certificate} & \(\eta\) en des points
; maximum sur \(D\) : grille fine puis L-BFGS-B borné depuis les 5
meilleurs maxima locaux \\
\texttt{joint\_descent} &
\(\min_{I\ge 0,\,z\in D}\tfrac12\lVert m-\sum_k I_k\varphi(z_k)\rVert^2+\lambda\sum_k I_k\)
par L-BFGS-B, gradient analytique \\
\texttt{debias} & moindres carrés positifs (NNLS) à positions fixées,
sans pénalité \\
\texttt{diagnostics} & \(\lVert r\rVert^2\) comparé à \(M\sigma^2\) \\
\texttt{match\_sources} & appariement hongrois ; erreurs de position et
d'intensité \\
\end{longtable}
}

\subsubsection{\texorpdfstring{2.3 LASSO positif sur grille
(\texttt{lasso.py})}{2.3 LASSO positif sur grille (lasso.py)}}\label{lasso-positif-sur-grille-lasso.py}

\begin{enumerate}
\def\labelenumi{\arabic{enumi}.}
\tightlist
\item
  \texttt{build\_grid(D,\ step)} construit la grille et la matrice
  \(\mathbf{G}\) (\(M\times P\)). Avec un pas de \(0{,}25\),
  \(P = 3321\) ; avec un pas de \(0{,}1\), \(P = 20\,301\).
\item
  \texttt{solve\_lasso} propose trois solveurs.

  \begin{itemize}
  \tightlist
  \item
    \textbf{FISTA} : pas \(1/L\) avec
    \(L = \lVert\mathbf{G}\rVert_2^2\), opérateur proximal
    \(c\mapsto\max(c-\lambda/L, 0)\), redémarrage adaptatif
    (O'Donoghue--Candès), arrêt sur les conditions KKT. On ajoute un
    \textbf{polissage} : LASSO exact restreint au support trouvé,
    accepté seulement s'il vérifie KKT sur toute la grille.
  \item
    \textbf{sklearn} :
    \texttt{Lasso(positive=True,\ fit\_intercept=False)} avec
    \(\alpha_{\mathrm{sklearn}} = \lambda/M\), car sklearn minimise
    \(\frac{1}{2M}\lVert m-\mathbf{G}c\rVert^2 + \alpha_{\mathrm{sklearn}}\lVert c\rVert_1\).
  \item
    \textbf{Ensemble actif} exact (algorithme du projet), qui sert de
    référence.
  \end{itemize}
\item
  \texttt{merge\_peaks} regroupe les nœuds actifs voisins (rayon
  \(1{,}5\) pas) : \(I = \sum_j c_j\) et \(z\) = barycentre pondéré.
\item
  \texttt{refine} : à \(K\) fixé, descente conjointe en \((I, z)\) sans
  grille.
\item
  \texttt{lasso\_pipeline} enchaîne grille, LASSO, regroupement,
  raffinement, test \(\max\eta\le 1\), débiaisage et diagnostics.
\end{enumerate}

\subsubsection{\texorpdfstring{2.4 Frank-Wolfe glissant
(\texttt{blasso.py})}{2.4 Frank-Wolfe glissant (blasso.py)}}\label{frank-wolfe-glissant-blasso.py}

Algorithme de Denoyelle, Duval, Peyré et Soubies (2019). On part de
\(\mu = 0\) (ou d'une initialisation donnée). À chaque itération :

\begin{enumerate}
\def\labelenumi{\arabic{enumi}.}
\tightlist
\item
  on calcule le résidu \(r = m - A(\mu)\) ;
\item
  on cherche \(z^{\star} = \arg\max_D \eta\) ;
\item
  on s'arrête si \(\eta(z^{\star})\le 1 + \mathrm{tol}\) : \(\mu\) est
  alors certifiée optimale ;
\item
  on ajoute un Dirac en \(z^{\star}\) ;
\item
  on réajuste les intensités par un LASSO positif exact sur les \(K\)
  colonnes ;
\item
  \textbf{glissement} : optimisation conjointe de \((I, z)\) par
  L-BFGS-B borné à \(D\) ;
\item
  on supprime les Diracs nuls et on fusionne ceux distants de moins de
  \(10^{-3}\) fois la diagonale de \(D\).
\end{enumerate}

La sortie contient l'historique (\(K\), \(\max\eta\), objectif) puis le
résultat débiaisé.

\subsubsection{\texorpdfstring{2.5 Hypothèses
(\texttt{hypotheses.py})}{2.5 Hypothèses (hypotheses.py)}}\label{hypothuxe8ses-hypotheses.py}

\texttt{check\_hypotheses(z,\ lam,\ sigma,\ ...)} :

\begin{enumerate}
\def\labelenumi{\arabic{enumi}.}
\tightlist
\item
  calcule \(\Gamma\), son rang (SVD) et son conditionnement ;
\item
  calcule \(p_V\) par pseudo-inverse et contrôle \(\eta_V(z_k) = 1\) et
  \(\nabla\eta_V(z_k) = 0\) (à \(10^{-15}\) près) ;
\item
  teste (H2) sur une grille fine hors de boules de rayon \(\rho\) autour
  des sources, puis raffine par L-BFGS-B et maximise sur chacun des 4
  côtés de \(D\) ;
\item
  teste (H3) avec les valeurs propres de la hessienne \(2\times 2\) de
  \(\eta_V\) ;
\item
  renvoie H1, H2, H3, \texttt{non\_degenerate}, le rang, le
  conditionnement, la marge et la \textbf{marge relative}, et trace la
  carte de \(\eta_V\) si demandé.
\end{enumerate}

\textbf{Choix numérique pour (H2).} Au bord d'une boule de rayon
\(\rho\), \(\eta_V\approx 1-\tfrac12\lvert\lambda_{\max}\rvert\rho^2\),
où \(\lambda_{\max}\) est la valeur propre de \(\nabla^2\eta_V(z_k)\) la
plus proche de 0. La marge \(1-\max\eta_V\) est donc petite par nature
(de \(10^{-7}\) à \(10^{-5}\) ici), et un seuil fixe donne des verdicts
arbitraires. On utilise :

\begin{itemize}
\tightlist
\item
  un seuil au niveau du bruit numérique,
  \(\max(10^{-9},\,100\times\text{erreur d'interpolation})\) ;
\item
  une \textbf{marge relative}
  \(\dfrac{1-\max\eta_V}{\tfrac12\lvert\lambda_{\max}\rvert\rho^2}\).
  Elle vaut environ 1 quand \(\eta_V\) décroît comme le prédit la
  hessienne, elle est \(\ll 1\) quand \(\eta_V\) est presque plat, et
  elle est négative quand (H2) échoue.
\end{itemize}

\texttt{hypotheses\_after\_fit(z\_hat,\ I\_hat)} applique la même
vérification aux positions \textbf{estimées} et renvoie
\texttt{reliable}.

\subsubsection{\texorpdfstring{2.6 Maximum de vraisemblance glouton et
choix de \(K\) par BIC
(\texttt{ml.py})}{2.6 Maximum de vraisemblance glouton et choix de K par BIC (ml.py)}}\label{maximum-de-vraisemblance-glouton-et-choix-de-k-par-bic-ml.py}

Cette méthode a été conçue pour le scénario du projet, où la grosse
source centrale fixe le niveau de bruit. Le BLASSO y a deux défauts :

\begin{enumerate}
\def\labelenumi{\arabic{enumi}.}
\tightlist
\item
  la pénalité \(\lambda\sum_k I_k\) biaise positions et intensités
  (biais \(O(\lambda)\)), alors qu'il faut \(\lambda\) de l'ordre de
  \(\lVert w\rVert\) pour éliminer les fausses sources ;
\item
  le nombre de sources n'est réglé qu'indirectement, par \(\lambda\).
\end{enumerate}

Le bruit étant gaussien de variance \(s_i^2\) connue, la
log-vraisemblance vaut, à une constante près,

\[
\log L(I, z) = -\frac12\sum_{i=1}^{M}\frac{\big(m_i - A(s)_i\big)^2}{s_i^2}
= -\frac12\,\big\lVert m_w - A_w(s)\big\rVert_2^2,
\qquad (m_w)_i = \frac{m_i}{s_i},\quad A_w(s)_i = \frac{A(s)_i}{s_i}.
\]

C'est le \textbf{blanchiment} : chaque capteur est divisé par l'écart
type de son bruit (\texttt{Setup.with\_weights(1/s)}). Pour le bruit du
projet, \(s_i = \sigma\,y_i\) (\texttt{relative\_noise\_std}), calculé
sur les mesures \textbf{brutes}, source centrale comprise.

\textbf{Algorithme \texttt{greedy\_ml}.} On part de \(K = 0\) et on
répète :

\begin{enumerate}
\def\labelenumi{\arabic{enumi}.}
\tightlist
\item
  \textbf{ajout} d'un Dirac au maximum de la corrélation
  \(z\mapsto\langle\psi_w(z), r_w\rangle\) entre le résidu blanchi et le
  noyau blanchi normalisé (\(\alpha = 1\)). C'est la même étape que dans
  Frank-Wolfe ;
\item
  intensités par moindres carrés positifs, puis \textbf{glissement sans
  pénalité} (\(\lambda = 0\)) de toutes les positions et intensités. On
  obtient ainsi l'estimateur du maximum de vraisemblance à \(K\) fixé
  (un minimum local) ;
\item
  calcul du \textbf{critère d'information}
\end{enumerate}

\[
C(K) = \big\lVert r_w\big\rVert_2^2 + \kappa\cdot 3K,
\qquad \kappa = \log M \ (\text{BIC}) \quad\text{ou}\quad \kappa = 2 \ (\text{AIC}),
\]

où \(3K\) est le nombre de paramètres (deux coordonnées et une intensité
par source).

On s'arrête quand \(C\) n'a pas diminué depuis deux ajouts, et on garde
le meilleur \(K\). Une \textbf{étape arrière} essaie ensuite de retirer
chaque source, tant que \(C\) diminue. Options : \texttt{n\_candidates}
(plusieurs positions candidates à chaque ajout),
\texttt{z\_init}/\texttt{I\_init} (départ imposé, par exemple la sortie
de Prony).

\subsubsection{\texorpdfstring{2.7 Borne de Cramér--Rao
(\texttt{ml.py})}{2.7 Borne de Cramér--Rao (ml.py)}}\label{borne-de-cramuxe9rrao-ml.py}

La borne de Cramér--Rao (CRB) donne l'erreur minimale de tout estimateur
sans biais. Avec \(\theta = (I_1,\dots,I_K, z_1,\dots,z_K)\) et

\[
\mathbf{D} = \Big[\varphi(z_1),\dots,\varphi(z_K)\;\Big\vert\;I_1\nabla\varphi(z_1),\dots,I_K\nabla\varphi(z_K)\Big]\in\mathbb{R}^{M\times 3K},
\]

la matrice de Fisher vaut
\(F = \mathbf{D}^{\top}\operatorname{diag}(s_i^{-2})\,\mathbf{D}\), et
pour tout estimateur sans biais

\[
\mathbb{E}\,\lVert\hat z_k - z_k\rVert^2 \;\ge\; \big(F^{-1}\big)_{z_k^{(1)} z_k^{(1)}} + \big(F^{-1}\big)_{z_k^{(2)} z_k^{(2)}} =: \mathrm{CRB}_k^2 .
\]

\texttt{cramer\_rao(z,\ I,\ s)} renvoie \(\mathrm{CRB}_k\) pour chaque
source. Dans les tableaux, « CRB » désigne \(\max_k \mathrm{CRB}_k\).
C'est une \textbf{limite physique} : aucun algorithme ne peut faire
mieux en moyenne.

\subsubsection{\texorpdfstring{2.8 Choix de
\(\lambda\)}{2.8 Choix de \textbackslash lambda}}\label{choix-de-lambda}

\begin{itemize}
\tightlist
\item
  Expériences génériques : \(\lambda = c\,\sigma\sqrt{M}\), avec
  \(c = 2\) par défaut (balayage \(c\in[0{,}25;\,16]\)).
\item
  Étude du noyau \(1/r^2\) :
  \(\lambda = c\,\sigma\sqrt{M}\,\lVert\Psi\rVert\) avec
  \(\lVert\Psi\rVert = \max_D\lVert\psi(z)\rVert\). Avec ce choix, le
  bruit seul ne peut pas faire dépasser \(1/c\) au certificat, quel que
  soit \(\alpha\).
\item
  Scénario du projet : \(\lambda = 1{,}5\,\sigma\lVert y\rVert\), une
  estimation de \(1{,}5\,\lVert w\rVert\).
\end{itemize}

\subsection{3. Structure du code}\label{structure-du-code}

\begin{Shaded}
\begin{Highlighting}[]
\NormalTok{sparse\_sources/}
\NormalTok{|{-}{-} \_\_init\_\_.py             API publique}
\NormalTok{|{-}{-} kernel.py               Setup, make\_setup, WeightedKernel, WhitenedKernel, phi et derivees}
\NormalTok{|{-}{-} common.py               forward, certificat, max\_certificate, descente, debias, diagnostics}
\NormalTok{|{-}{-} lasso.py                build\_grid, FISTA / sklearn / ensemble actif, merge\_peaks, refine}
\NormalTok{|{-}{-} blasso.py               sliding\_frank\_wolfe}
\NormalTok{|{-}{-} hypotheses.py           check\_hypotheses, hypotheses\_after\_fit, format\_report}
\NormalTok{|{-}{-} ml.py                   greedy\_ml (maximum de vraisemblance + BIC), cramer\_rao}
\NormalTok{|{-}{-} plotting.py             style des figures}
\NormalTok{|{-}{-} experiments.py          experiences generales (3 geometries, comparaison Prony)}
\NormalTok{|{-}{-} experiments\_1r2.py      etude ciblee du noyau 1/r\^{}2}
\NormalTok{|{-}{-} experiments\_projet.py   scenario du projet (source centrale), Cramer{-}Rao}
\NormalTok{|{-}{-} tests/                  70 tests pytest}
\NormalTok{|{-}{-} results/                CSV (versionnes) et PNG (ignores par .gitignore)}
\NormalTok{|{-}{-} README.md               notice courte}
\NormalTok{|{-}{-} COMPTE\_RENDU.md         ce document}
\NormalTok{\textasciigrave{}{-}{-} COMPTE\_RENDU.tex        ce document converti en LaTeX par Pandoc}
\end{Highlighting}
\end{Shaded}

\texttt{sparse\_sources} importe \texttt{particle\_sources} (noyau,
\texttt{Box}, \texttt{CircleSensors}, configuration, Prony, génération).
\textbf{Aucun fichier existant du projet n'a été modifié.} Toutes les
fonctions acceptent un argument facultatif \texttt{setup=} ; sans lui,
elles utilisent la géométrie de \texttt{config/}.

\subsection{4. Validation}\label{validation}

\texttt{pytest\ sparse\_sources/tests} exécute \textbf{70 tests} en
environ 30 s. Les 67 tests existants du projet passent toujours.

{\def\LTcaptype{none} % do not increment counter
\begin{longtable}[]{@{}
  >{\raggedright\arraybackslash}p{(\linewidth - 2\tabcolsep) * \real{0.5000}}
  >{\raggedright\arraybackslash}p{(\linewidth - 2\tabcolsep) * \real{0.5000}}@{}}
\toprule\noalign{}
\begin{minipage}[b]{\linewidth}\raggedright
Ce qui est testé
\end{minipage} & \begin{minipage}[b]{\linewidth}\raggedright
Résultat
\end{minipage} \\
\midrule\noalign{}
\endhead
\bottomrule\noalign{}
\endlastfoot
\(\nabla\varphi\), \(\nabla^2\varphi\) (brut, \(h = 20\), pondéré,
blanchi) contre différences finies & erreur \(\le 10^{-6}\) (en pratique
\(10^{-10}\)) \\
\(\varphi\) identique au noyau du projet (\texttt{intensity\_at}) &
oui \\
\(\eta_V(z_k) = 1\) et \(\nabla\eta_V(z_k) = 0\) & à \(10^{-10}\)
près \\
Une source bien vue (\(\alpha = 1\), ou \(h = 20\)) & H1, H2, H3
vraies \\
Noyau brut \(1/r^2\) dans le plan, une source & H2 et H3 fausses \\
Accord avec \texttt{duval\_peyre.diagnose} du projet & mêmes verdicts \\
FISTA, sklearn, ensemble actif & même objectif (\(10^{-4}\) ;
\(10^{-3}\) pour sklearn) \\
Frank-Wolfe sans bruit & positions à moins de \(10^{-4}\) \\
Bruit \(10^{-4}\), \(10^{-3}\), \(10^{-2}\) (Frank-Wolfe et LASSO) & bon
\(K\), erreur \(\le 300\,\sigma_{\mathrm{rel}}\) \\
Erreur proportionnelle au bruit (\(10^{-5}\) à \(10^{-3}\)) & rapport
constant à un facteur 3 près \\
\texttt{greedy\_ml}, 2 sources, bruit de la source centrale & bon \(K\),
erreur \(< 5\times\) CRB, \(\chi^2\) compatible avec le bruit \\
\texttt{greedy\_ml}, bruit seul & au plus 1 source retenue \\
Efficacité : erreur rms du ML pour une source (30 tirages) & entre
\(0{,}6\) et \(1{,}5\) fois la CRB \\
\end{longtable}
}

\subsection{5. Résultats}\label{ruxe9sultats}

\subsubsection{\texorpdfstring{5.1 Trois géométries, trois
configurations
(\texttt{experiments.py})}{5.1 Trois géométries, trois configurations (experiments.py)}}\label{trois-guxe9omuxe9tries-trois-configurations-experiments.py}

Géométries : \textbf{plan\_brut} (projet : \(h = 0\), \(\alpha = 0\)),
\textbf{plan\_normalise} (\(h = 0\), \(\alpha = 1\)),
\textbf{hauteur\_20} (\(h = 20\), \(\alpha = 0\)). Configurations : 3
sources séparées ; 2 sources à distance 1 ; 2 sources avec des capteurs
sur un cercle 2 fois plus grand (« loin »).

\textbf{Hypothèses} (\texttt{hypotheses.csv}, cartes
\texttt{certificat\_*.png}) :

{\def\LTcaptype{none} % do not increment counter
\begin{longtable}[]{@{}
  >{\raggedright\arraybackslash}p{(\linewidth - 14\tabcolsep) * \real{0.1250}}
  >{\raggedright\arraybackslash}p{(\linewidth - 14\tabcolsep) * \real{0.1250}}
  >{\raggedright\arraybackslash}p{(\linewidth - 14\tabcolsep) * \real{0.1250}}
  >{\raggedright\arraybackslash}p{(\linewidth - 14\tabcolsep) * \real{0.1250}}
  >{\raggedright\arraybackslash}p{(\linewidth - 14\tabcolsep) * \real{0.1250}}
  >{\raggedright\arraybackslash}p{(\linewidth - 14\tabcolsep) * \real{0.1250}}
  >{\raggedright\arraybackslash}p{(\linewidth - 14\tabcolsep) * \real{0.1250}}
  >{\raggedright\arraybackslash}p{(\linewidth - 14\tabcolsep) * \real{0.1250}}@{}}
\toprule\noalign{}
\begin{minipage}[b]{\linewidth}\raggedright
géométrie
\end{minipage} & \begin{minipage}[b]{\linewidth}\raggedright
config
\end{minipage} & \begin{minipage}[b]{\linewidth}\raggedright
H1
\end{minipage} & \begin{minipage}[b]{\linewidth}\raggedright
H2
\end{minipage} & \begin{minipage}[b]{\linewidth}\raggedright
H3
\end{minipage} & \begin{minipage}[b]{\linewidth}\raggedright
marge
\end{minipage} & \begin{minipage}[b]{\linewidth}\raggedright
marge relative
\end{minipage} & \begin{minipage}[b]{\linewidth}\raggedright
cond(\(\Gamma\)) normalisé
\end{minipage} \\
\midrule\noalign{}
\endhead
\bottomrule\noalign{}
\endlastfoot
plan\_brut & séparées & oui & non & non & \(-41{,}8\) & \(<0\) & 122 \\
plan\_brut & proches & oui & non & non & \(-13{,}7\) & \(<0\) & 3 980 \\
plan\_brut & loin & oui & non & non & \(-0{,}50\) & \(<0\) & 94 \\
plan\_normalise & séparées & oui & oui & oui & \(8{,}5\times 10^{-6}\) &
1,24 & 92 \\
plan\_normalise & proches & oui & oui & oui & \(1{,}1\times 10^{-6}\) &
0,70 & 3 950 \\
plan\_normalise & loin & oui & oui & oui & \(6{,}5\times 10^{-6}\) &
1,32 & 85 \\
hauteur\_20 & séparées & oui & oui & oui & \(8{,}4\times 10^{-7}\) &
1,42 & \textbf{127 000} \\
hauteur\_20 & proches & oui & oui & oui & \(1{,}5\times 10^{-7}\) & 0,70
& \textbf{293 000} \\
hauteur\_20 & loin & oui & non & non & \(-0{,}012\) & \(<0\) & 507 \\
\end{longtable}
}

\textbf{Comparaison des méthodes} (\(c = 2\), 4 tirages). Taux de
réussite, défini comme « bon \(K\) et erreur de position inférieure au
quart de la séparation » :

{\def\LTcaptype{none} % do not increment counter
\begin{longtable}[]{@{}
  >{\raggedright\arraybackslash}p{(\linewidth - 8\tabcolsep) * \real{0.2000}}
  >{\raggedright\arraybackslash}p{(\linewidth - 8\tabcolsep) * \real{0.2000}}
  >{\raggedright\arraybackslash}p{(\linewidth - 8\tabcolsep) * \real{0.2000}}
  >{\raggedright\arraybackslash}p{(\linewidth - 8\tabcolsep) * \real{0.2000}}
  >{\raggedright\arraybackslash}p{(\linewidth - 8\tabcolsep) * \real{0.2000}}@{}}
\toprule\noalign{}
\begin{minipage}[b]{\linewidth}\raggedright
géométrie, config
\end{minipage} & \begin{minipage}[b]{\linewidth}\raggedright
\(\sigma/\mathrm{rms}\)
\end{minipage} & \begin{minipage}[b]{\linewidth}\raggedright
LASSO
\end{minipage} & \begin{minipage}[b]{\linewidth}\raggedright
LASSO + raffinement
\end{minipage} & \begin{minipage}[b]{\linewidth}\raggedright
Frank-Wolfe
\end{minipage} \\
\midrule\noalign{}
\endhead
\bottomrule\noalign{}
\endlastfoot
plan\_brut, toutes & \(10^{-4}\) à \(10^{-2}\) & 0 \% (23 à 40 sources)
& 0 \% & 0 \% \\
plan\_normalise, séparées & \(10^{-4}\) / \(10^{-3}\) / \(10^{-2}\) & 0
/ 0 / 25 \% & 50 / 75 / 75 \% & \textbf{100 / 100 / 75 \%} \\
plan\_normalise, proches & \(10^{-4}\) / \(10^{-3}\) / \(10^{-2}\) & 50
/ 25 / 0 \% & 100 / 25 / 0 \% & \textbf{100 / 25 / 0 \%} \\
plan\_normalise, loin & \(10^{-4}\) / \(10^{-3}\) / \(10^{-2}\) & 50 /
50 / 75 \% & 50 / 100 / 25 \% & \textbf{100 / 100 / 25 \%} \\
hauteur\_20, séparées et proches & \(10^{-4}\) à \(10^{-2}\) & 0 \% (une
seule source) & 0 \% & 0 \% \\
\end{longtable}
}

Le LASSO seul surestime souvent \(K\), à cause de petits pics voisins
non fusionnés ; le raffinement corrige la plupart des cas, et le
Frank-Wolfe est le plus fiable.

\textbf{Erreur en fonction du bruit} (\texttt{erreur\_sigma.png},
Frank-Wolfe, plan\_normalise) : l'erreur de position est
\textbf{proportionnelle à \(\sigma\)} de \(10^{-5}\) à \(10^{-2}\), avec
un rapport \(\text{err}/\sigma_{\mathrm{rel}}\) constant : environ 72
(séparées), 97 (loin) et 330 (proches). C'est la loi
\(O(\lVert w\rVert)\) du théorème ; la constante grandit quand les
sources se rapprochent.

\textbf{Balayage de \(\lambda\)} (\texttt{balayage\_lambda.png}) : avec
\(\alpha = 1\), \(K\) est exact pour \(c\) entre \(0{,}5\) et \(4\)
environ. Un \(c\) trop petit ajoute des sources parasites, un \(c\) trop
grand fusionne les sources proches. Avec le noyau brut, on trouve 20 à
35 sources de trop pour tout \(c\).

\subsubsection{\texorpdfstring{5.2 Étude ciblée du noyau \(1/r^2\)
(\texttt{experiments\_1r2.py})}{5.2 Étude ciblée du noyau 1/r\^{}2 (experiments\_1r2.py)}}\label{uxe9tude-cibluxe9e-du-noyau-1r2-experiments_1r2.py}

Sources tirées comme dans \texttt{config/} (positions uniformes,
intensités uniformes sur le simplexe de somme 100), 20 tirages par
\(K\), \(\sigma = 10^{-3}\times\mathrm{rms}\),
\(\lambda = 2\sigma\sqrt{M}\lVert\Psi\rVert\).

\textbf{Effet de la pondération \(\alpha\)} (\texttt{r2\_alpha.png}) :

{\def\LTcaptype{none} % do not increment counter
\begin{longtable}[]{@{}
  >{\raggedright\arraybackslash}p{(\linewidth - 6\tabcolsep) * \real{0.2500}}
  >{\raggedright\arraybackslash}p{(\linewidth - 6\tabcolsep) * \real{0.2500}}
  >{\raggedright\arraybackslash}p{(\linewidth - 6\tabcolsep) * \real{0.2500}}
  >{\raggedright\arraybackslash}p{(\linewidth - 6\tabcolsep) * \real{0.2500}}@{}}
\toprule\noalign{}
\begin{minipage}[b]{\linewidth}\raggedright
\(\alpha\)
\end{minipage} & \begin{minipage}[b]{\linewidth}\raggedright
condition vraie (\(K\) = 1 / 2 / 3 / 4 / 5)
\end{minipage} & \begin{minipage}[b]{\linewidth}\raggedright
réussite Frank-Wolfe (\(K\) = 1 / 2 / 3 / 4 / 5)
\end{minipage} & \begin{minipage}[b]{\linewidth}\raggedright
temps médian
\end{minipage} \\
\midrule\noalign{}
\endhead
\bottomrule\noalign{}
\endlastfoot
0 (brut) & 0 / 0 / 0 / 0 / 0 \% & 0 / 0 / 0 / 0 / 0 \% & 0,59 s (non
convergé) \\
0,5 & 0 / 0 / 0 / 0 / 0 \% & 0 / 0 / 0 / 0 / 0 \% & 1,29 s \\
0,75 & 35 / 25 / 0 / 0 / 0 \% & 35 / 25 / 0 / 0 / 0 \% & 0,21 s \\
\textbf{1} & 100 / 100 / 20 / 20 / 0 \% & \textbf{100 / 95 / 25 / 15 / 5
\%} & \textbf{0,05 s} \\
\textbf{1,25} & 100 / 100 / 25 / 25 / 0 \% & \textbf{100 / 95 / 30 / 15
/ 0 \%} & 0,04 s \\
1,5 & 100 / 100 / 35 / 30 / 0 \% & 100 / 95 / 25 / 5 / 0 \% & 0,04 s \\
2 & 100 / 90 / 50 / 30 / 0 \% & 100 / 80 / 20 / 5 / 0 \% & 0,04 s \\
\end{longtable}
}

Il faut \(\alpha\ge 1\) pour une seule source. Pour \(\alpha = 1\) et
une source en \(z_0\), sans bruit et pour \(\lambda\to 0\), le
certificat vaut

\[
\eta(z) = \big\langle\psi(z), \psi(z_0)\big\rangle \le \lVert\psi(z)\rVert\,\lVert\psi(z_0)\rVert = 1
\]

par Cauchy--Schwarz, avec égalité seulement si \(\psi(z) = \psi(z_0)\),
donc en \(z = z_0\). Au-delà de \(\alpha\approx 1{,}25\), la condition
est plus souvent vraie, mais avec un conditionnement plus mauvais, et la
réussite baisse. \textbf{Recommandation : \(\alpha = 1\) (ou
\(1{,}25\)).}

\textbf{Pouvoir prédictif de \texttt{check\_hypotheses}}
(\texttt{r2\_prediction.png}, 100 tirages par \(\alpha\)) :

{\def\LTcaptype{none} % do not increment counter
\begin{longtable}[]{@{}llll@{}}
\toprule\noalign{}
\(\alpha\) & condition vraie & réussite si vraie & réussite si fausse \\
\midrule\noalign{}
\endhead
\bottomrule\noalign{}
\endlastfoot
0,75 & 12 / 100 & 92 \% & 1 \% \\
\textbf{1} & 48 / 100 & \textbf{94 \%} & \textbf{6 \%} \\
1,25 & 50 / 100 & 94 \% & 2 \% \\
1,5 & 53 / 100 & 83 \% & 2 \% \\
2 & 54 / 100 & 74 \% & 2 \% \\
\end{longtable}
}

Le signe de la marge relative sépare presque parfaitement les réussites
des échecs : la vérification des hypothèses est un \textbf{outil de
diagnostic fiable}, utilisable a posteriori avec
\texttt{hypotheses\_after\_fit}.

\textbf{Limite de résolution} (\texttt{r2\_resolution.png}, paire
verticale au centre, \(\alpha = 1\), 5 tirages) :

{\def\LTcaptype{none} % do not increment counter
\begin{longtable}[]{@{}
  >{\raggedright\arraybackslash}p{(\linewidth - 4\tabcolsep) * \real{0.3333}}
  >{\raggedright\arraybackslash}p{(\linewidth - 4\tabcolsep) * \real{0.3333}}
  >{\raggedright\arraybackslash}p{(\linewidth - 4\tabcolsep) * \real{0.3333}}@{}}
\toprule\noalign{}
\begin{minipage}[b]{\linewidth}\raggedright
\(\sigma/\mathrm{rms}\)
\end{minipage} & \begin{minipage}[b]{\linewidth}\raggedright
séparation minimale résolue (100 \%)
\end{minipage} & \begin{minipage}[b]{\linewidth}\raggedright
erreur à \(d = 2\)
\end{minipage} \\
\midrule\noalign{}
\endhead
\bottomrule\noalign{}
\endlastfoot
\(10^{-5}\) & 0,5 & \(1{,}3\times 10^{-3}\) \\
\(10^{-4}\) & 1 & \(1{,}3\times 10^{-2}\) \\
\(10^{-3}\) & 1,5 & 0,14 \\
\(10^{-2}\) & 4 & (\(K\) faux à \(d = 2\)) \\
\end{longtable}
}

On observe environ \(d^{\star}\propto\sigma^{1/4}\) : le bruit
admissible décroît comme \(d^4\), ce qui est typique de la
super-résolution de sources positives proches. À \(d = 0{,}25\), la
condition est vraie mais
\(\operatorname{cond}(\Gamma) = 2{,}5\times 10^5\), et aucun niveau de
bruit testé ne permet la reconstruction.

\textbf{Distance aux capteurs} (paire horizontale de séparation 2,
\(\alpha = 1\), \(\sigma = 10^{-3}\)) :

{\def\LTcaptype{none} % do not increment counter
\begin{longtable}[]{@{}
  >{\raggedright\arraybackslash}p{(\linewidth - 6\tabcolsep) * \real{0.2500}}
  >{\raggedright\arraybackslash}p{(\linewidth - 6\tabcolsep) * \real{0.2500}}
  >{\raggedright\arraybackslash}p{(\linewidth - 6\tabcolsep) * \real{0.2500}}
  >{\raggedright\arraybackslash}p{(\linewidth - 6\tabcolsep) * \real{0.2500}}@{}}
\toprule\noalign{}
\begin{minipage}[b]{\linewidth}\raggedright
\(y\) (centre = 10)
\end{minipage} & \begin{minipage}[b]{\linewidth}\raggedright
distance au capteur le plus proche
\end{minipage} & \begin{minipage}[b]{\linewidth}\raggedright
condition
\end{minipage} & \begin{minipage}[b]{\linewidth}\raggedright
réussite
\end{minipage} \\
\midrule\noalign{}
\endhead
\bottomrule\noalign{}
\endlastfoot
10 à 14 & 10,2 à 7,1 & vraie & 100 \% \\
16 & 5,1 & vraie & 60 \% \\
17,5 / 18,5 / 19,25 & 3,6 / 2,6 / 1,9 & \textbf{fausse} & \textbf{0
\%} \\
\end{longtable}
}

Une \textbf{source seule} reste reconstructible près des capteurs. C'est
la \textbf{paire} qui échoue : près d'un capteur, \(\psi(z)\) est dominé
par ce capteur, et deux sources à la même distance des mêmes capteurs
ont des signatures presque identiques
(\texttt{r2\_certificat\_pres\_capteurs.png}).

\textbf{Nombre de capteurs} (\texttt{r2\_capteurs.png}) : passer de
\(M = 16\) à \(M = 256\) change peu de chose. Les échecs sont
\textbf{géométriques} ; ajouter des capteurs sur le \textbf{même cercle}
n'apporte pas d'information nouvelle.

\subsubsection{\texorpdfstring{5.3 La condition est qualitative : le cas
\(h = 20\)}{5.3 La condition est qualitative : le cas h = 20}}\label{la-condition-est-qualitative-le-cas-h-20}

Avec des capteurs à \(h = 20\) au-dessus du plan, \(G(x,\cdot)\) n'est
plus sous-harmonique sur \(D\), car

\[
\Delta_z G = \frac{4\,(r^2-h^2)}{(r^2+h^2)^3} < 0 \quad\text{pour } r < h,
\]

et la condition est \textbf{vraie} pour les configurations séparées et
proches. Pourtant \textbf{la reconstruction échoue même sans bruit} : le
noyau est si lisse que les colonnes sont presque colinéaires, avec
\(\operatorname{cond}(\Gamma)\approx 10^5\). \textbf{À retenir :
(H1)--(H3) ne suffisent pas en pratique ; il faut aussi regarder
\(\operatorname{cond}(\Gamma)\) et la marge relative.}

\subsubsection{\texorpdfstring{5.4 Première comparaison avec
\texttt{reconstruct.py}
(Prony)}{5.4 Première comparaison avec reconstruct.py (Prony)}}\label{premiuxe8re-comparaison-avec-reconstruct.py-prony}

Scénario identique à \texttt{scripts/reconstruct.py} : \(n\) sources
tirées selon \texttt{config/}, source centrale d'intensité 5000 mesurée
puis retirée, bruit relatif, 8 tirages par case, mêmes mesures pour
toutes les méthodes. Taux de réussite (et erreur de position médiane) :

{\def\LTcaptype{none} % do not increment counter
\begin{longtable}[]{@{}
  >{\raggedright\arraybackslash}p{(\linewidth - 12\tabcolsep) * \real{0.1429}}
  >{\raggedright\arraybackslash}p{(\linewidth - 12\tabcolsep) * \real{0.1429}}
  >{\raggedright\arraybackslash}p{(\linewidth - 12\tabcolsep) * \real{0.1429}}
  >{\raggedright\arraybackslash}p{(\linewidth - 12\tabcolsep) * \real{0.1429}}
  >{\raggedright\arraybackslash}p{(\linewidth - 12\tabcolsep) * \real{0.1429}}
  >{\raggedright\arraybackslash}p{(\linewidth - 12\tabcolsep) * \real{0.1429}}
  >{\raggedright\arraybackslash}p{(\linewidth - 12\tabcolsep) * \real{0.1429}}@{}}
\toprule\noalign{}
\begin{minipage}[b]{\linewidth}\raggedright
\(n\)
\end{minipage} & \begin{minipage}[b]{\linewidth}\raggedright
\(\sigma\)
\end{minipage} & \begin{minipage}[b]{\linewidth}\raggedright
Prony (config)
\end{minipage} & \begin{minipage}[b]{\linewidth}\raggedright
Prony (\(5\sigma\))
\end{minipage} & \begin{minipage}[b]{\linewidth}\raggedright
Frank-Wolfe brut
\end{minipage} & \begin{minipage}[b]{\linewidth}\raggedright
LASSO normalisé
\end{minipage} & \begin{minipage}[b]{\linewidth}\raggedright
Frank-Wolfe normalisé
\end{minipage} \\
\midrule\noalign{}
\endhead
\bottomrule\noalign{}
\endlastfoot
2 & \(10^{-4}\) & 88 \% (0,02) & 25 \% & 0 \% & \textbf{100 \%} (0,05) &
88 \% (0,05) \\
2 & \(10^{-3}\) & 25 \% (0,41) & 25 \% & 0 \% & \textbf{75 \%} (0,50) &
62 \% (0,50) \\
2 & \(10^{-2}\) & 0 \% & 12 \% & 0 \% & 12 \% & 12 \% \\
3 & \(10^{-4}\) & \textbf{62 \%} (0,16) & 25 \% & 0 \% & 25 \% & 25
\% \\
3 & \(10^{-3}\) & 25 \% & 25 \% & 0 \% & 0 \% & 0 \% \\
5 & \(10^{-4}\) & 12 \% & \textbf{25 \%} & 0 \% & 0 \% & 0 \% \\
\end{longtable}
}

Prony est meilleur à très faible bruit et pour plus de sources ; le
BLASSO pondéré est plus robuste à bruit moyen pour 2 sources. Cette
comparaison a motivé la recherche d'une meilleure méthode (section
suivante).

\subsubsection{\texorpdfstring{5.5 Scénario du projet : maximum de
vraisemblance et borne de Cramér--Rao
(\texttt{experiments\_projet.py})}{5.5 Scénario du projet : maximum de vraisemblance et borne de Cramér--Rao (experiments\_projet.py)}}\label{scuxe9nario-du-projet-maximum-de-vraisemblance-et-borne-de-cramuxe9rrao-experiments_projet.py}

Données identiques à \texttt{scripts/reconstruct.py} : \(n\) sources
tirées selon \texttt{config/} (intensités uniformes sur le simplexe de
somme 100, positions uniformes), source centrale \(I_c = 5000\) ajoutée
puis retirée, bruit relatif \(\sigma\). 20 tirages par case, mêmes
mesures pour toutes les méthodes :

\begin{itemize}
\tightlist
\item
  \texttt{prony} : \texttt{cfg.invert}, exactement comme
  \texttt{reconstruct.py} ;
\item
  \texttt{sfw} : Frank-Wolfe glissant, BLASSO pondéré \(\alpha = 1\),
  \(\lambda = 1{,}5\,\sigma\lVert y\rVert\) (meilleure méthode des
  sections précédentes) ;
\item
  \texttt{ml\_bic} : maximum de vraisemblance glouton, bruit blanchi,
  \(K\) par BIC (section 2.6) ;
\item
  \texttt{ml\_bic\_multi} : idem avec 3 candidats par ajout et un second
  départ depuis Prony ; on garde le meilleur critère.
\end{itemize}

Pour chaque tirage, on calcule aussi la borne de Cramér--Rao (CRB). Un
tirage est dit \textbf{faisable} si
\(2\,\mathrm{CRB} < \tfrac14\min(d_{\min}, 4)\), où \(d_{\min}\) est la
plus petite distance entre deux sources : une méthode optimale a alors
de bonnes chances de passer le critère de réussite.

\textbf{Taux de réussite} (\texttt{projet\_methodes\_resume.csv},
\texttt{projet\_methodes.png}) :

{\def\LTcaptype{none} % do not increment counter
\begin{longtable}[]{@{}
  >{\raggedright\arraybackslash}p{(\linewidth - 14\tabcolsep) * \real{0.1250}}
  >{\raggedright\arraybackslash}p{(\linewidth - 14\tabcolsep) * \real{0.1250}}
  >{\raggedright\arraybackslash}p{(\linewidth - 14\tabcolsep) * \real{0.1250}}
  >{\raggedright\arraybackslash}p{(\linewidth - 14\tabcolsep) * \real{0.1250}}
  >{\raggedright\arraybackslash}p{(\linewidth - 14\tabcolsep) * \real{0.1250}}
  >{\raggedright\arraybackslash}p{(\linewidth - 14\tabcolsep) * \real{0.1250}}
  >{\raggedright\arraybackslash}p{(\linewidth - 14\tabcolsep) * \real{0.1250}}
  >{\raggedright\arraybackslash}p{(\linewidth - 14\tabcolsep) * \real{0.1250}}@{}}
\toprule\noalign{}
\begin{minipage}[b]{\linewidth}\raggedright
\(n\)
\end{minipage} & \begin{minipage}[b]{\linewidth}\raggedright
\(\sigma\)
\end{minipage} & \begin{minipage}[b]{\linewidth}\raggedright
Prony
\end{minipage} & \begin{minipage}[b]{\linewidth}\raggedright
Frank-Wolfe
\end{minipage} & \begin{minipage}[b]{\linewidth}\raggedright
\textbf{ML + BIC}
\end{minipage} & \begin{minipage}[b]{\linewidth}\raggedright
ML multi
\end{minipage} & \begin{minipage}[b]{\linewidth}\raggedright
tirages faisables
\end{minipage} & \begin{minipage}[b]{\linewidth}\raggedright
ML + BIC sur les faisables
\end{minipage} \\
\midrule\noalign{}
\endhead
\bottomrule\noalign{}
\endlastfoot
1 & \(10^{-4}\) & 100 \% & 100 \% & \textbf{100 \%} & 100 \% & 100 \% &
100 \% \\
1 & \(10^{-3}\) & 15 \% & 100 \% & \textbf{100 \%} & 100 \% & 100 \% &
100 \% \\
1 & \(10^{-2}\) & 0 \% & 100 \% & \textbf{100 \%} & 100 \% & 85 \% & 100
\% \\
2 & \(10^{-4}\) & 85 \% & 90 \% & \textbf{95 \%} & 95 \% & 85 \% & 94
\% \\
2 & \(3\times 10^{-4}\) & 40 \% & 75 \% & \textbf{85 \%} & 85 \% & 80 \%
& 94 \% \\
2 & \(10^{-3}\) & 15 \% & 70 \% & \textbf{75 \%} & 75 \% & 80 \% & 94
\% \\
2 & \(3\times 10^{-3}\) & 0 \% & 40 \% & \textbf{75 \%} & 75 \% & 50 \%
& 90 \% \\
2 & \(10^{-2}\) & 5 \% & 15 \% & \textbf{30 \%} & 30 \% & 20 \% & 100
\% \\
3 & \(10^{-4}\) & 55 \% & 40 \% & \textbf{80 \%} & 75 \% & 75 \% & 100
\% \\
3 & \(3\times 10^{-4}\) & 25 \% & 20 \% & \textbf{65 \%} & 65 \% & 35 \%
& 100 \% \\
3 & \(10^{-3}\) & 10 \% & 5 \% & \textbf{40 \%} & 35 \% & 25 \% & 100
\% \\
3 & \(10^{-2}\) & 0 \% & 0 \% & 0 \% & 0 \% & 0 \% & sans objet \\
5 & \(10^{-4}\) & 5 \% & 0 \% & \textbf{50 \%} & 50 \% & 20 \% & 100
\% \\
5 & \(3\times 10^{-4}\) & 5 \% & 0 \% & \textbf{15 \%} & 15 \% & 5 \% &
100 \% \\
5 & \(10^{-3}\) & 0 \% & 0 \% & 5 \% & 5 \% & 5 \% & 100 \% \\
\end{longtable}
}

\textbf{Erreur de position médiane comparée à la CRB} (même fichier) :

{\def\LTcaptype{none} % do not increment counter
\begin{longtable}[]{@{}llllll@{}}
\toprule\noalign{}
\(n\) & \(\sigma\) & Prony & Frank-Wolfe & ML + BIC & CRB médiane \\
\midrule\noalign{}
\endhead
\bottomrule\noalign{}
\endlastfoot
1 & \(10^{-3}\) & 0,059 & 0,023 & \textbf{0,022} & 0,025 \\
1 & \(10^{-2}\) & 0,84 & 0,23 & \textbf{0,23} & 0,25 \\
2 & \(10^{-4}\) & 0,014 & 0,035 & \textbf{0,014} & 0,016 \\
2 & \(10^{-3}\) & 0,32 & 0,29 & \textbf{0,13} & 0,16 \\
2 & \(3\times 10^{-3}\) & 1,6 & 0,66 & \textbf{0,36} & 0,49 \\
3 & \(10^{-4}\) & 0,14 & 0,72 & \textbf{0,15} & 0,22 \\
3 & \(10^{-3}\) & 4,3 & 0,93 & \textbf{0,42} & 2,2 \\
5 & \(10^{-4}\) & 0,88 & 1,6 & \textbf{0,60} & 1,04 \\
\end{longtable}
}

\textbf{Lecture.}

\begin{enumerate}
\def\labelenumi{\arabic{enumi}.}
\tightlist
\item
  \textbf{Le maximum de vraisemblance + BIC est la meilleure méthode
  dans toutes les cases}, ou à égalité avec la meilleure. Le gain est
  net dès 2 sources : 75 \% contre 15 \% pour Prony à \(n = 2\),
  \(\sigma = 10^{-3}\) ; 80 \% contre 55 \% à \(n = 3\),
  \(\sigma = 10^{-4}\) ; 50 \% contre 5 \% à \(n = 5\),
  \(\sigma = 10^{-4}\).
\item
  \textbf{Son erreur atteint la borne de Cramér--Rao.} Pour 1 et 2
  sources, l'erreur médiane vaut \(0{,}7\) à \(0{,}9\) fois la CRB (la
  CRB porte sur la pire source, l'erreur médiane sur un tirage typique).
  La méthode est donc \textbf{efficace} au sens statistique : aucun
  estimateur sans biais ne peut faire mieux. Le test
  \texttt{test\_cramer\_rao\_scaling\_and\_efficiency} le vérifie pour
  une source.
\item
  \textbf{Quand le problème est faisable, la méthode réussit} : 90 à 100
  \% de réussite sur les tirages faisables, dans toutes les cases. Ses
  échecs correspondent presque tous à des tirages où la CRB dépasse la
  tolérance : aucune méthode ne pourrait les réussir de manière fiable.
\item
  \textbf{D'où vient le gain par rapport au BLASSO} : surtout de la
  suppression de la pénalité lors du glissement (plus de biais
  \(O(\lambda)\) sur les positions) et du choix de \(K\) par BIC. Le
  blanchiment seul apporte peu ici : la source centrale, au centre du
  cercle, donne presque le même bruit à tous les capteurs (dans le
  prototype, le Frank-Wolfe blanchi et le non blanchi donnent des
  résultats très proches).
\item
  \textbf{Prony} est excellent pour une source à très faible bruit, mais
  il surestime \(K\) dès que \(\sigma\ge 10^{-3}\) (4 à 7 sources
  estimées), car son seuil de rang n'est pas adapté au bruit.
\item
  \textbf{Les variantes « multi »} (plusieurs candidats, départ depuis
  Prony) n'apportent rien : elles trouvent parfois un critère plus bas,
  mais une reconstruction moins bonne. La méthode simple suffit.
\item
  Quand la CRB est très grande (par exemple \(n = 3\),
  \(\sigma = 10^{-3}\) : CRB \(= 2{,}2\)), l'erreur des méthodes peut
  être inférieure à la CRB. Elles sont alors biaisées (régularisation
  implicite), et la borne indique surtout que la position de la pire
  source n'est pas déterminée par les données.
\end{enumerate}

\textbf{Choix du nombre de sources : BIC contre AIC}
(\texttt{projet\_critere.csv}) :

{\def\LTcaptype{none} % do not increment counter
\begin{longtable}[]{@{}
  >{\raggedright\arraybackslash}p{(\linewidth - 10\tabcolsep) * \real{0.1667}}
  >{\raggedright\arraybackslash}p{(\linewidth - 10\tabcolsep) * \real{0.1667}}
  >{\raggedright\arraybackslash}p{(\linewidth - 10\tabcolsep) * \real{0.1667}}
  >{\raggedright\arraybackslash}p{(\linewidth - 10\tabcolsep) * \real{0.1667}}
  >{\raggedright\arraybackslash}p{(\linewidth - 10\tabcolsep) * \real{0.1667}}
  >{\raggedright\arraybackslash}p{(\linewidth - 10\tabcolsep) * \real{0.1667}}@{}}
\toprule\noalign{}
\begin{minipage}[b]{\linewidth}\raggedright
\(n\)
\end{minipage} & \begin{minipage}[b]{\linewidth}\raggedright
\(\sigma\)
\end{minipage} & \begin{minipage}[b]{\linewidth}\raggedright
réussite BIC
\end{minipage} & \begin{minipage}[b]{\linewidth}\raggedright
réussite AIC
\end{minipage} & \begin{minipage}[b]{\linewidth}\raggedright
trop de sources (BIC / AIC)
\end{minipage} & \begin{minipage}[b]{\linewidth}\raggedright
trop peu (BIC / AIC)
\end{minipage} \\
\midrule\noalign{}
\endhead
\bottomrule\noalign{}
\endlastfoot
2 & \(10^{-4}\) & \textbf{95 \%} & 85 \% & 5 / 15 \% & 0 / 0 \% \\
2 & \(10^{-3}\) & \textbf{75 \%} & 70 \% & 5 / 15 \% & 15 / 10 \% \\
3 & \(10^{-4}\) & \textbf{80 \%} & 60 \% & 0 / 20 \% & 0 / 0 \% \\
3 & \(10^{-3}\) & \textbf{40 \%} & 35 \% & 0 / 10 \% & 50 / 25 \% \\
5 & \(10^{-4}\) & \textbf{50 \%} & 40 \% & 0 / 20 \% & 35 / 25 \% \\
\end{longtable}
}

Le BIC (\(\kappa = \log M\approx 4{,}2\)) est meilleur partout ; l'AIC
(\(\kappa = 2\)) ajoute des sources parasites. Les erreurs du BIC sont
surtout des sources \textbf{manquantes} à fort bruit : des sources
faibles ou proches que les données ne permettent pas de voir.

\textbf{Effet de la source centrale}
(\texttt{projet\_centrale\_resume.csv}, \texttt{projet\_centrale.png},
\(n = 3\), \(\sigma = 10^{-3}\), 20 tirages ; le facteur vaut
\(I_c/100\)) :

{\def\LTcaptype{none} % do not increment counter
\begin{longtable}[]{@{}
  >{\raggedright\arraybackslash}p{(\linewidth - 10\tabcolsep) * \real{0.1667}}
  >{\raggedright\arraybackslash}p{(\linewidth - 10\tabcolsep) * \real{0.1667}}
  >{\raggedright\arraybackslash}p{(\linewidth - 10\tabcolsep) * \real{0.1667}}
  >{\raggedright\arraybackslash}p{(\linewidth - 10\tabcolsep) * \real{0.1667}}
  >{\raggedright\arraybackslash}p{(\linewidth - 10\tabcolsep) * \real{0.1667}}
  >{\raggedright\arraybackslash}p{(\linewidth - 10\tabcolsep) * \real{0.1667}}@{}}
\toprule\noalign{}
\begin{minipage}[b]{\linewidth}\raggedright
facteur
\end{minipage} & \begin{minipage}[b]{\linewidth}\raggedright
Prony
\end{minipage} & \begin{minipage}[b]{\linewidth}\raggedright
Frank-Wolfe
\end{minipage} & \begin{minipage}[b]{\linewidth}\raggedright
ML + BIC
\end{minipage} & \begin{minipage}[b]{\linewidth}\raggedright
erreur ML + BIC
\end{minipage} & \begin{minipage}[b]{\linewidth}\raggedright
CRB médiane
\end{minipage} \\
\midrule\noalign{}
\endhead
\bottomrule\noalign{}
\endlastfoot
0 (pas de source centrale) & 55 \% & 35 \% & \textbf{100 \%} & 0,034 &
0,052 \\
5 & 55 \% & 35 \% & \textbf{70 \%} & 0,22 & 0,30 \\
50 (projet) & 10 \% & 5 \% & \textbf{40 \%} & 0,42 & 2,2 \\
500 & 0 \% & 0 \% & 0 \% & 1,2 & 21 \\
\end{longtable}
}

\textbf{C'est le résultat principal.} Sans source centrale, le maximum
de vraisemblance réussit dans 100 \% des cas. Avec la source centrale du
projet, la CRB est multipliée par environ 42. Le bruit de la source
centrale, d'écart type \(\sigma I_c/R^2 = 40\,\sigma\) par capteur,
reste dans les données après soustraction. Il domine le signal des
petites sources, qui vaut environ 7 en moyenne quadratique : avec
\(\sigma = 10^{-3}\), le rapport signal sur bruit par capteur n'est que
d'environ 175.

\textbf{Configuration par défaut du projet}
(\(N\sim\mathrm{Poisson}(20)\), \(\sigma = 1\,\%\), 10 tirages,
\texttt{projet\_defaut.csv}) :

\begin{itemize}
\tightlist
\item
  les tirages contiennent 12 à 29 sources, d'intensité maximale 11 à 26
  ;
\item
  la CRB conjointe médiane est comprise entre \(10^4\) et
  \(5\times 10^5\) : les positions ne sont \textbf{pas déterminées} par
  les données ;
\item
  même en supposant toutes les autres sources connues, une source n'a
  une CRB inférieure à \(0{,}5\) que dans 0 à 4 cas par tirage ;
\item
  ML + BIC retient 1 à 5 sources, Frank-Wolfe 1 à 4, Prony 0 à 9.
\end{itemize}

Dans cette configuration, on ne peut espérer retrouver que les quelques
sources les plus intenses et les mieux placées. Il faut soit beaucoup
moins de bruit, soit beaucoup moins de sources, soit une autre géométrie
de capteurs (section 8).

\subsection{6. Où ça ne marche pas, et
pourquoi}\label{ouxf9-uxe7a-ne-marche-pas-et-pourquoi}

\subsubsection{\texorpdfstring{6.1 BLASSO standard avec le noyau
\(1/r^2\) : échec
systématique}{6.1 BLASSO standard avec le noyau 1/r\^{}2 : échec systématique}}\label{blasso-standard-avec-le-noyau-1r2-uxe9chec-systuxe9matique}

\begin{itemize}
\tightlist
\item
  \textbf{Constat} : la condition est fausse dans 100 \% des cas, et la
  réussite est nulle. On obtient 20 à 40 sources parasites collées au
  bord de \(D\), surtout aux coins.
\item
  \textbf{Cause 1 : sous-harmonicité.} En dimension 2,
  \(\Delta_z\big(1/\lVert x-z\rVert^2\big) = 4/\lVert x-z\rVert^4 > 0\)
  : chaque \(G(x_i,\cdot)\) est sous-harmonique. Lorsque \(p_V\ge 0\),
  \(\eta_V\) est sous-harmonique et ne peut pas avoir de maximum strict
  à l'intérieur de \(D\) : la trace de \(\nabla^2\eta_V(z_k)\) est
  positive (H3 fausse) et le maximum est au bord (H2 fausse,
  \(\max\eta_V\) entre 17 et 43 près des coins).
\item
  \textbf{Cause 2 : biais de la pénalité.} \(\lVert\varphi(z)\rVert\)
  varie d'un facteur 215 sur \(D\). Le BLASSO « paie » moins cher pour
  remplacer une source centrale par plusieurs sources de bord ; la
  solution du LASSO sur grille est d'ailleurs un anneau de sources de
  bord, même sans bruit.
\item
  \textbf{Remède} : la pondération \(\alpha\approx 1\) (section 1.2).
\end{itemize}

\subsubsection{6.2 Sources proches}\label{sources-proches}

La séparation minimale résolue va de \(0{,}5\) à \(4\) selon le bruit
(\(d^{\star}\propto\sigma^{1/4}\)). Le conditionnement de \(\Gamma\)
explose (environ 4 000 à \(d = 1\) et 250 000 à \(d = 0{,}25\)).

\subsubsection{6.3 Zone aveugle près des
capteurs}\label{zone-aveugle-pruxe8s-des-capteurs}

Le cercle des capteurs passe par les coins de \(D\). Les \textbf{paires}
situées à moins de 4 à 5 unités d'un capteur ne sont plus séparables.
Comme les positions uniformes tombent souvent près du bord, cela
explique une grande partie des échecs pour \(K\ge 3\).

\subsubsection{6.4 Bruit de la source
centrale}\label{bruit-de-la-source-centrale}

C'est la limite dominante dans le scénario du projet (section 5.5).
Retirer la source centrale enlève son signal, mais pas son bruit,
d'écart type \(40\,\sigma\) par capteur. Aucune méthode ne peut
descendre sous la borne de Cramér--Rao, qui croît proportionnellement à
ce bruit.

\subsubsection{6.5 Nombre de sources}\label{nombre-de-sources}

Avec des positions aléatoires, la probabilité d'avoir une paire proche
ou une paire près des capteurs croît vite avec \(K\). Dans la
configuration par défaut (\(K\approx 20\), \(\sigma = 1\,\%\)), la borne
de Cramér--Rao conjointe dépasse \(10^4\) : le problème est hors de
portée de toute méthode.

\subsubsection{\texorpdfstring{6.6 Noyau trop lisse (capteurs hors du
plan, \(h\)
grand)}{6.6 Noyau trop lisse (capteurs hors du plan, h grand)}}\label{noyau-trop-lisse-capteurs-hors-du-plan-h-grand}

La condition est vraie mais \(\operatorname{cond}(\Gamma)\approx 10^5\)
: c'est inexploitable numériquement (section 5.3).

\subsubsection{6.7 Problèmes numériques
rencontrés}\label{probluxe8mes-numuxe9riques-rencontruxe9s}

\begin{itemize}
\tightlist
\item
  \textbf{FISTA est lent} sur une grille fine, car les colonnes voisines
  sont presque colinéaires : 27 s au pas de \(0{,}1\) contre \(0{,}17\)
  s pour l'ensemble actif. Le polissage sur le support corrige la
  non-convergence.
\item
  \textbf{sklearn} s'arrête sur son nombre maximal d'itérations (écart
  d'objectif d'environ \(3\times 10^{-4}\)).
\item
  \textbf{LASSO + raffinement} : un petit pic parasite peut rester
  bloqué dans un minimum local (intensité \(\approx 0\), donc gradient
  en position \(\approx 0\)).
\item
  \textbf{\(\lambda\) très petit} (moins de \(10^{-6}\lambda_{\max}\)) :
  le critère \(\eta\le 1+10^{-6}\) n'est plus atteignable numériquement.
\item
  \textbf{Seuil de (H2)} : un seuil fixe donne des verdicts arbitraires,
  d'où la marge relative (section 2.5).
\item
  \textbf{Maximum de vraisemblance} : le problème à \(K\) fixé n'est pas
  convexe. Un meilleur critère ne garantit pas une meilleure
  reconstruction : partir de Prony ou essayer plusieurs candidats trouve
  parfois un optimum de critère plus bas mais plus loin de la vérité
  (section 5.5).
\end{itemize}

\subsection{7. Temps de calcul}\label{temps-de-calcul}

Machine de développement (Windows, Python 3.14, NumPy 2.3, SciPy 1.16),
un cœur, valeurs médianes.

{\def\LTcaptype{none} % do not increment counter
\begin{longtable}[]{@{}
  >{\raggedright\arraybackslash}p{(\linewidth - 4\tabcolsep) * \real{0.3333}}
  >{\raggedright\arraybackslash}p{(\linewidth - 4\tabcolsep) * \real{0.3333}}
  >{\raggedright\arraybackslash}p{(\linewidth - 4\tabcolsep) * \real{0.3333}}@{}}
\toprule\noalign{}
\begin{minipage}[b]{\linewidth}\raggedright
Brique
\end{minipage} & \begin{minipage}[b]{\linewidth}\raggedright
Paramètres
\end{minipage} & \begin{minipage}[b]{\linewidth}\raggedright
Temps
\end{minipage} \\
\midrule\noalign{}
\endhead
\bottomrule\noalign{}
\endlastfoot
\texttt{make\_grid} (matrice \(\mathbf{G}\)) & pas 0,1 (\(P = 20\,301\))
& 0,07 s \\
Frank-Wolfe glissant, \(\alpha = 1\) & \(K\) = 1 / 2 / 3 / 5 / 8 & 0,004
/ 0,017 / 0,045 / 0,10 / 0,44 s \\
Frank-Wolfe glissant, noyau brut & plafonné à 10 itérations & 0,4 à 1,3
s \\
\texttt{lasso\_pipeline}, ensemble actif, \(K = 3\) & pas 0,5 / 0,25 /
0,1 & 0,03 / 0,04 / 0,17 s \\
\texttt{lasso\_pipeline}, FISTA, \(K = 3\) & pas 0,5 / 0,25 / 0,1 & 0,86
/ 2,0 / 27 s \\
\texttt{lasso\_pipeline}, sklearn, \(K = 3\) & pas 0,5 / 0,25 / 0,1 &
1,5 / 1,9 / 16 s \\
\texttt{check\_hypotheses} & grille 0,05, \(K\) = 1 à 8 & 0,25 à 0,30
s \\
Prony (\texttt{prony\_circle}) & \(M = 64\) & environ 0,001 s \\
\texttt{greedy\_ml} (ML + BIC) & \(n\) = 1 / 2 / 3 / 5,
\(\sigma = 10^{-4}\) & 0,06 / 0,12 / 0,19 / 0,42 s \\
\texttt{greedy\_ml}, 3 candidats + départ Prony & \(n\) = 1 / 3 / 5 &
0,24 / 0,65 / 1,4 s \\
\texttt{cramer\_rao} & \(K\le 30\) & moins de 0,01 s \\
\texttt{pytest\ sparse\_sources/tests} & 70 tests & environ 30 s \\
\texttt{experiments.py} & tout & 336 s \\
\texttt{experiments\_1r2.py} & tout (20 tirages par \(K\)) & 391 s \\
\texttt{experiments\_projet.py} & tout (20 tirages par case) & 366 s \\
\end{longtable}
}

Le Frank-Wolfe glissant est le plus rapide quand la condition tient
(environ \(K+1\) itérations). Le maximum de vraisemblance glouton est 3
à 10 fois plus lent : il ajoute en général deux sources de trop avant de
s'arrêter, puis fait l'étape arrière. Il reste sous la demi-seconde pour
5 sources.

\subsection{8. Pistes d'amélioration}\label{pistes-damuxe9lioration}

\textbf{Modélisation}

\begin{enumerate}
\def\labelenumi{\arabic{enumi}.}
\tightlist
\item
  \textbf{Utiliser le maximum de vraisemblance glouton + BIC}
  (\texttt{greedy\_ml}) comme méthode par défaut dans le scénario du
  projet, et le BLASSO pondéré (\(\alpha = 1\)) quand on veut un
  certificat d'optimalité.
\item
  \textbf{Réduire le bruit de la source centrale} plutôt que chercher un
  meilleur algorithme : mesurer plus longtemps, protéger les capteurs de
  la source centrale (collimation), ou estimer son intensité
  conjointement si elle n'est pas connue exactement.
\item
  \textbf{Bruit de Poisson} (\texttt{poisson\_counts}) : remplacer
  l'attache quadratique par la log-vraisemblance de Poisson (divergence
  de Kullback--Leibler) ; la même méthode gloutonne s'applique.
\end{enumerate}

\textbf{Géométrie des capteurs} (facteur limitant principal après le
bruit)

\begin{enumerate}
\def\labelenumi{\arabic{enumi}.}
\setcounter{enumi}{3}
\tightlist
\item
  \textbf{Éloigner les capteurs des coins de \(D\)} (cercle plus grand,
  \texttt{radius\_factor}) pour supprimer la zone aveugle. La borne de
  Cramér--Rao et \texttt{check\_hypotheses} permettent de comparer des
  géométries avant de construire quoi que ce soit.
\item
  \textbf{Diversifier les capteurs} plutôt que les multiplier sur le
  même cercle (plusieurs cercles, capteurs à l'intérieur, hauteur
  modérée \(h\sim 2\) à \(5\)).
\end{enumerate}

\textbf{Algorithmes}

\begin{enumerate}
\def\labelenumi{\arabic{enumi}.}
\setcounter{enumi}{5}
\tightlist
\item
  \textbf{Choix de \(K\)} : le BIC est un bon compromis ; un test du
  \(\chi^2\) sur \(\lVert r_w\rVert^2\) ou une validation croisée sur
  les capteurs sont des alternatives.
\item
  \textbf{LASSO sur grille} : remplacer FISTA par l'ensemble actif ou
  par une descente par coordonnées avec élimination (\emph{screening}).
\end{enumerate}

\textbf{Théorie}

\begin{enumerate}
\def\labelenumi{\arabic{enumi}.}
\setcounter{enumi}{7}
\tightlist
\item
  \textbf{Séparation dans la métrique du noyau} (métrique de
  Fisher--Rao, Poon--Keriven--Peyré) pour expliquer quantitativement la
  zone aveugle et la loi \(d^{\star}\propto\sigma^{1/4}\).
\item
  \textbf{Constantes quantitatives} : relier \(\lambda_0\) et \(c_0\) à
  la marge relative et à \(\operatorname{cond}(\Gamma)\).
\end{enumerate}

\subsection{9. Écarts par rapport à la demande
initiale}\label{uxe9carts-par-rapport-uxe0-la-demande-initiale}

\begin{itemize}
\tightlist
\item
  Toutes les fonctions acceptent un argument supplémentaire
  \texttt{setup=} (géométrie, \(h\), \(\alpha\), poids des capteurs).
\item
  \texttt{debias(z,\ I\_init,\ m)} prend aussi les mesures \(m\),
  nécessaires au NNLS.
\item
  \texttt{check\_hypotheses} a des paramètres supplémentaires et renvoie
  en plus la marge relative et des indicateurs de bruit.
\item
  Ajouts non demandés, jugés nécessaires : la pondération \(\alpha\), le
  solveur exact à ensemble actif, le polissage de FISTA, la marge
  relative de (H2), l'étude ciblée du noyau \(1/r^2\), la comparaison
  avec Prony, le maximum de vraisemblance glouton, le blanchiment et la
  borne de Cramér--Rao.
\item
  Les \textbf{PNG} ne sont pas versionnés (le \texttt{.gitignore} du
  dépôt ignore \texttt{*.png}) ; les scripts d'expériences les
  régénèrent. Les CSV sont versionnés.
\end{itemize}

\subsection{10. Reproduire}\label{reproduire}

\begin{Shaded}
\begin{Highlighting}[]
\ExtensionTok{pip}\NormalTok{ install }\AttributeTok{{-}e} \StringTok{".[dev]"}\NormalTok{ scikit{-}learn}
\ExtensionTok{pytest}\NormalTok{ sparse\_sources/tests                   }\CommentTok{\# 70 tests, environ 30 s}
\ExtensionTok{python}\NormalTok{ sparse\_sources/experiments.py          }\CommentTok{\# environ 6 min}
\ExtensionTok{python}\NormalTok{ sparse\_sources/experiments\_1r2.py      }\CommentTok{\# environ 6,5 min}
\ExtensionTok{python}\NormalTok{ sparse\_sources/experiments\_projet.py   }\CommentTok{\# scenario du projet}
\end{Highlighting}
\end{Shaded}

Toutes les graines sont fixées : les résultats sont reproductibles.

Pour obtenir une version LaTeX ou PDF du compte rendu :

\begin{Shaded}
\begin{Highlighting}[]
\ExtensionTok{pandoc}\NormalTok{ COMPTE\_RENDU.md }\AttributeTok{{-}s} \AttributeTok{{-}o}\NormalTok{ COMPTE\_RENDU.tex}
\ExtensionTok{pandoc}\NormalTok{ COMPTE\_RENDU.md }\AttributeTok{{-}s} \AttributeTok{{-}{-}pdf{-}engine}\OperatorTok{=}\NormalTok{xelatex }\AttributeTok{{-}o}\NormalTok{ COMPTE\_RENDU.pdf}
\end{Highlighting}
\end{Shaded}

\subsection{11. Conclusion : quelle méthode utiliser
?}\label{conclusion-quelle-muxe9thode-utiliser}

Taux de réussite sur le scénario du projet (section 5.5, source centrale
comprise, 20 tirages par case) :

{\def\LTcaptype{none} % do not increment counter
\begin{longtable}[]{@{}
  >{\raggedright\arraybackslash}p{(\linewidth - 8\tabcolsep) * \real{0.2000}}
  >{\raggedright\arraybackslash}p{(\linewidth - 8\tabcolsep) * \real{0.2000}}
  >{\raggedright\arraybackslash}p{(\linewidth - 8\tabcolsep) * \real{0.2000}}
  >{\raggedright\arraybackslash}p{(\linewidth - 8\tabcolsep) * \real{0.2000}}
  >{\raggedright\arraybackslash}p{(\linewidth - 8\tabcolsep) * \real{0.2000}}@{}}
\toprule\noalign{}
\begin{minipage}[b]{\linewidth}\raggedright
sources
\end{minipage} & \begin{minipage}[b]{\linewidth}\raggedright
bruit \(\sigma\)
\end{minipage} & \begin{minipage}[b]{\linewidth}\raggedright
Prony
\end{minipage} & \begin{minipage}[b]{\linewidth}\raggedright
Frank-Wolfe (\(\alpha = 1\))
\end{minipage} & \begin{minipage}[b]{\linewidth}\raggedright
\textbf{ML + BIC}
\end{minipage} \\
\midrule\noalign{}
\endhead
\bottomrule\noalign{}
\endlastfoot
1 & \(10^{-4}\) (faible) & 100 \% & 100 \% & \textbf{100 \%} \\
1 & \(10^{-3}\) & 15 \% & 100 \% & \textbf{100 \%} \\
2 & \(10^{-4}\) (faible) & 85 \% & 90 \% & \textbf{95 \%} \\
2 & \(3\times 10^{-3}\) & 0 \% & 40 \% & \textbf{75 \%} \\
3 & \(10^{-4}\) (faible) & 55 \% & 40 \% & \textbf{80 \%} \\
3 & \(10^{-3}\) & 10 \% & 5 \% & \textbf{40 \%} \\
5 & \(10^{-4}\) (faible) & 5 \% & 0 \% & \textbf{50 \%} \\
\end{longtable}
}

\begin{itemize}
\tightlist
\item
  \textbf{Prony est bon quand le bruit est très faible.} C'est une
  méthode algébrique, exacte sans bruit. Mais dès que \(\sigma\) atteint
  \(10^{-3}\), il trouve trop de sources (4 à 7) et s'effondre, même
  avec une seule source (15 \%).
\item
  \textbf{Frank-Wolfe résiste bien au bruit, mais seulement pour 1 ou 2
  sources.} À partir de 3 sources, il bute sur la condition (H2) :
  sources proches, sources près des capteurs.
\item
  \textbf{Avec beaucoup de sources, les deux sont mauvais.} Prony fait
  un peu mieux à faible bruit (55 \% contre 40 \% pour 3 sources), mais
  avec 5 sources ils sont tous les deux proches de 0 \%. Prony n'est
  donc pas « bon pour beaucoup de sources ».
\item
  \textbf{ML + BIC fait au moins aussi bien que la meilleure des deux
  dans chaque case}, et nettement mieux dès 2 sources avec du bruit.
\end{itemize}

\textbf{Tableau de conclusion :}

{\def\LTcaptype{none} % do not increment counter
\begin{longtable}[]{@{}
  >{\raggedright\arraybackslash}p{(\linewidth - 6\tabcolsep) * \real{0.2500}}
  >{\raggedright\arraybackslash}p{(\linewidth - 6\tabcolsep) * \real{0.2500}}
  >{\raggedright\arraybackslash}p{(\linewidth - 6\tabcolsep) * \real{0.2500}}
  >{\raggedright\arraybackslash}p{(\linewidth - 6\tabcolsep) * \real{0.2500}}@{}}
\toprule\noalign{}
\begin{minipage}[b]{\linewidth}\raggedright
situation
\end{minipage} & \begin{minipage}[b]{\linewidth}\raggedright
meilleure méthode classique
\end{minipage} & \begin{minipage}[b]{\linewidth}\raggedright
méthode à utiliser
\end{minipage} & \begin{minipage}[b]{\linewidth}\raggedright
ce qui limite
\end{minipage} \\
\midrule\noalign{}
\endhead
\bottomrule\noalign{}
\endlastfoot
bruit très faible & Prony & \textbf{ML + BIC} & rien en pratique (1 à 3
sources) \\
bruit moyen, 1 ou 2 sources & Frank-Wolfe (\(\alpha = 1\)) & \textbf{ML
+ BIC} & le bruit, au niveau de la borne de Cramér--Rao \\
3 sources ou plus & aucune & \textbf{ML + BIC} & la condition (H2) pour
le BLASSO, le bruit pour le ML \\
bruit fort, source centrale du projet & aucune & \textbf{ML + BIC} & la
borne de Cramér--Rao : le problème n'est plus résoluble \\
\end{longtable}
}

En pratique, la méthode à retenir est \texttt{greedy\_ml} (maximum de
vraisemblance + BIC). Son erreur atteint la borne de Cramér--Rao : pour
faire mieux, il faut changer les données (moins de bruit sur la source
centrale, autre géométrie de capteurs), pas l'algorithme. Prony et
Frank-Wolfe ne servent plus que de points de comparaison.

\subsection{Références}\label{ruxe9fuxe9rences}

\begin{itemize}
\tightlist
\item
  V. Duval, G. Peyré, \emph{Exact support recovery for sparse spikes
  deconvolution}, Foundations of Computational Mathematics, 2015.
\item
  Q. Denoyelle, V. Duval, G. Peyré, E. Soubies, \emph{The sliding
  Frank-Wolfe algorithm and its application to super-resolution
  microscopy}, Inverse Problems, 2019.
\item
  Q. Denoyelle, V. Duval, G. Peyré, \emph{Support recovery for sparse
  super-resolution of positive measures}, Journal of Fourier Analysis
  and Applications, 2017.
\item
  C. Poon, N. Keriven, G. Peyré, \emph{The geometry of off-the-grid
  compressed sensing}, Foundations of Computational Mathematics, 2023.
\item
  A. Beck, M. Teboulle, \emph{A fast iterative shrinkage-thresholding
  algorithm for linear inverse problems}, SIAM Journal on Imaging
  Sciences, 2009.
\item
  B. O'Donoghue, E. Candès, \emph{Adaptive restart for accelerated
  gradient schemes}, Foundations of Computational Mathematics, 2015.
\item
  G. Schwarz, \emph{Estimating the dimension of a model}, Annals of
  Statistics, 1978 (critère BIC).
\item
  S. M. Kay, \emph{Fundamentals of Statistical Signal Processing:
  Estimation Theory}, Prentice Hall, 1993 (borne de Cramér--Rao).
\end{itemize}

\end{document}
